\documentclass[11pt]{article}
\pagestyle{plain}

% \pgfplotsset{compat=1.18}
\usepackage[margin=1in]{geometry}
\usepackage{amsmath,amsfonts,amssymb,amsthm}
\usepackage[ruled,vlined,linesnumbered,commentsnumbered]{algorithm2e}
\SetKwInput{KwIn}{Input}
\SetKwInput{KwOut}{Output}
\SetKwInput{KwGlobal}{Global variable}
\SetKwInput{KwOracle}{Oracle}
\SetKwComment{Comment}{// }{}
\usepackage{graphicx}
\usepackage{enumerate}
\usepackage{bbm}
\usepackage{verbatim}
\usepackage{hyperref,color}
\usepackage[capitalize,nameinlink]{cleveref}
\usepackage[dvipsnames]{xcolor} 
\hypersetup{
	colorlinks=true,
	pdfpagemode=UseNone,
	citecolor=OliveGreen,
	linkcolor=NavyBlue,
	urlcolor=Magenta,
	pdfstartview=FitW
}
\usepackage{appendix}
\usepackage{makecell}
\usepackage{tablefootnote}
\crefname{appsec}{Appendix}{Appendices}
\usepackage{tikz}
\usepackage{pgfplots}
\usepackage{xifthen}
\usepackage{subcaption}
\usepackage{hhline}

\theoremstyle{plain}
\newtheorem{theorem}{Theorem}[section]
\newtheorem{proposition}[theorem]{Proposition}
\newtheorem{lemma}[theorem]{Lemma}
\newtheorem{corollary}[theorem]{Corollary}
\newtheorem{conjecture}[theorem]{Conjecture}
\newtheorem{problem}[theorem]{Problem}
\newtheorem{claim}[theorem]{Claim}
\newtheorem{fact}[theorem]{Fact}

\theoremstyle{definition}
\newtheorem{definition}[theorem]{Definition}
\newtheorem{example}[theorem]{Example}
\newtheorem{observation}[theorem]{Observation}
\newtheorem{assumption}[theorem]{Assumption}
\newtheorem*{assumption*}{Assumption}
\newtheorem{condition}[theorem]{Condition}

\theoremstyle{remark}
\newtheorem{remark}[theorem]{Remark}

\crefname{lemma}{Lemma}{Lemmas}
\crefname{theorem}{Theorem}{Theorems}
\crefname{definition}{Definition}{Definitions}
\crefname{fact}{Fact}{Facts}
\crefname{claim}{Claim}{Claims}
\crefname{proposition}{Proposition}{Propositions}


\newcommand{\dif}{\,\mathrm{d}}
%\newcommand{\E}{\mathbb{E}}
%\newcommand{\Var}{\mathrm{Var}}
\newcommand{\Cov}{\mathrm{Cov}}
\newcommand{\MI}{\mathrm{I}}
%\newcommand{\Ent}{\mathrm{Ent}}
\newcommand{\ind}{\mathbbm{1}}
\newcommand{\allone}{\mathbf{1}}
\newcommand{\tr}{\mathrm{tr}}
\newcommand{\wt}{\mathrm{weight}}
\DeclareMathOperator*{\argmax}{arg\,max}
\newcommand{\norm}[1]{\left\lVert #1 \right\rVert}
\newcommand{\TV}[2]{\left\|#1 - #2\right\|_{\mathrm{TV}}}
\newcommand{\ceil}[1]{\left\lceil #1 \right\rceil}
\newcommand{\floor}[1]{\left\lfloor #1 \right\rfloor}
\newcommand{\ip}[2]{\left\langle #1 , #2 \right\rangle}
\newcommand{\grad}{\nabla}
\newcommand{\hessian}{\nabla^2}
\newcommand{\trans}{\intercal}
\newcommand{\diag}{\mathrm{diag}}
\newcommand{\poly}{\mathrm{poly}}
\newcommand{\dist}{\mathrm{dist}}
\newcommand{\sgn}{\mathrm{sgn}}
\newcommand{\fpras}{\mathsf{FPRAS}}
\newcommand{\fpaus}{\mathsf{FPAUS}}
\newcommand{\fptas}{\mathsf{FPTAS}}

\newcommand{\eps}{\varepsilon}

\newcommand{\N}{\mathbb{N}}
\newcommand{\Z}{\mathbb{Z}}
\newcommand{\Q}{\mathbb{Q}}
\newcommand{\R}{\mathbb{R}}
\newcommand{\C}{\mathbb{C}}

\renewcommand{\AA}{\mathcal{A}}
\newcommand{\BB}{\mathcal{B}}
\newcommand{\CC}{\mathcal{C}}
\newcommand{\DD}{\mathcal{D}}
\newcommand{\EE}{\mathcal{E}}
\newcommand{\FF}{\mathcal{F}}
\newcommand{\GG}{\mathcal{G}}
\newcommand{\HH}{\mathcal{H}}
\newcommand{\II}{\mathcal{I}}
\newcommand{\JJ}{\mathcal{J}}
\newcommand{\KK}{\mathcal{K}}
\newcommand{\LL}{\mathcal{L}}
\newcommand{\MM}{\mathcal{M}}
\newcommand{\NN}{\mathcal{N}}
\newcommand{\OO}{\mathcal{O}}
\newcommand{\PP}{\mathcal{P}}
\newcommand{\QQ}{\mathcal{Q}}
\newcommand{\RR}{\mathcal{R}}
\renewcommand{\SS}{\mathcal{S}}
\newcommand{\TT}{\mathcal{T}}
\newcommand{\UU}{\mathcal{U}}
\newcommand{\VV}{\mathcal{V}}
\newcommand{\WW}{\mathcal{W}}
\newcommand{\XX}{\mathcal{X}}
\newcommand{\YY}{\mathcal{Y}}
\newcommand{\ZZ}{\mathcal{Z}}

\newcommand{\KL}[2]{D_{\mathrm{KL}} \left( #1 \,\Vert\, #2 \right)}
\newcommand{\KLbig}[2]{D_{\mathrm{KL}} \left( #1 \,\big\Vert\, #2 \right)}
\newcommand{\KLBig}[2]{D_{\mathrm{KL}} \left( #1 \,\Big\Vert\, #2 \right)}

\newcommand{\SKL}[2]{D_{\mathrm{SKL}} \left( #1 \,\Vert\, #2 \right)}
\newcommand{\SKLbig}[2]{D_{\mathrm{SKL}} \left( #1 \,\big\Vert\, #2 \right)}
\newcommand{\SKLBig}[2]{D_{\mathrm{SKL}} \left( #1 \,\Big\Vert\, #2 \right)}

\newcommand{\qandq}{\quad\text{and}\quad}
\newcommand{\supp}{\mathrm{supp}}

\newcommand{\DKL}[2]{\-D_{\-{KL}}\left(#1 \parallel #2\right)}
\newcommand{\DTV}[2]{\-D_{\mathrm{TV}}\left({#1},{#2}\right)}
% \newcommand{\dist}{\mathrm{dist}}
\newcommand{\e}{\mathrm{e}}
\renewcommand{\epsilon}{\varepsilon}
\renewcommand{\emptyset}{\varnothing}
% \newcommand{\norm}[1]{\left\Vert#1\right\Vert}
\newcommand{\set}[1]{\left\{#1\right\}}
\newcommand{\tuple}[1]{\left(#1\right)} 
% \newcommand{\eps}{\varepsilon}
\newcommand{\inner}[2]{\left\langle #1,#2\right\rangle}
\newcommand{\tp}{\tuple}
\newcommand{\ol}{\overline}
\newcommand{\abs}[1]{\left\vert#1\right\vert}
\newcommand{\ctp}[1]{\left\lceil#1\right\rceil}
\newcommand{\ftp}[1]{\left\lfloor#1\right\rfloor}
\newcommand{\Ind}[1]{\-{Ind}_{#1}}
\newcommand{\todo}[1]{{\color{red} TODO: #1}}

\newcommand{\Cor}{\Psi^{\mathrm{sym}}}
\newcommand{\cor}{\Psi^{\-{cor}}}
\newcommand{\Inf}{\Psi}
% \newcommand{\diag}{\-{diag}}

\def\*#1{\boldsymbol{#1}} % Use \*A for \mathbf{A}
\def\+#1{\mathcal{#1}} % Use \+A for \mathcal{A}
\def\-#1{\mathrm{#1}} % Use \-A for \mathrm{A}
\def\=#1{\mathbb{#1}} % Use \^A for \mathbb{A}
\def\!#1{\mathfrak{#1}} % Use \!A for \mathfrak{A}

\newcommand*\diff{\mathop{}\!\mathrm{d}}

% probability
% \DeclareMathOperator*{\oPr}{\mathbf{Pr}}
\def\oPr{\mathop{\mathrm{Pr}}}
\renewcommand{\Pr}[2][]{ \ifthenelse{\isempty{#1}}
  {\oPr\left[#2\right]}
  {\oPr_{#1}\left[#2\right]} } % Use \Pr[a]{b} for \mathbf{Pr}_a[b], \Pr{b} for  \mathbf{Pr}[b]

% expectation: relies on the package "dsfont"
% \DeclareMathOperator*{\oE}{\mathbb{E}}
\def\oE{\mathop{\mathbb{E}}}
\newcommand{\E}[2][]{ \ifthenelse{\isempty{#1}}
  {\oE\left[#2\right]}
  {\oE_{#1}\left[#2\right]} }

% variance
%\DeclareMathOperator*{\oVar}{\mathbf{Var}}
\def\oVar{\mathrm{Var}}
\newcommand{\Var}[2][]{ \ifthenelse{\isempty{#1}}
  {\oVar\left[#2\right]}
  {\oVar_{#1}\left[#2\right]} }

% entropy
% \DeclareMathOperator*{\oEnt}{\mathbf{Ent}}
\def\oEnt{\mathrm{Ent}}
\newcommand{\Ent}[2][]{ \ifthenelse{\isempty{#1}}
  {\oEnt\left[#2\right]}
  {\oEnt_{#1}\left[#2\right]} }

\newcommand{\PhiEnt}[2][]{ \ifthenelse{\isempty{#1}}
  {\oEnt^\Psi\left[#2\right]}
  {\oEnt^\Psi_{#1}\left[#2\right]} }


\newenvironment{Exampleparagraph}[1][]
{
    \par 
    \noindent 
    \textbf{\color{blue} #1}
    \par 
    \vspace{0.5em} 
    \itshape 
}
{
    \par 
    \vspace{1em} 
}

%\input{local-macro}

\newcommand{\RED}[1]{\textcolor{red}{#1}}
\newcommand{\zc}[1]{\textcolor{red}{ZC: #1}}

\newcommand{\bern}[1]{\mathrm{Bern}\left(#1\right)}
\newcommand{\pll}{\parallel}



\title{Gaussian Isopermetric Profile Entropic Independence and Contraction}
\date{\today}

\pgfplotsset{compat=1.18} 
\begin{document}

\maketitle
\section{Setup}
\begin{definition}[$\Psi$-entropy]
    Given convex function $\Psi: [0,\infty)\mapsto \R$, define $\Psi$-entropy $\PhiEnt{X}$ as follows
    $$
    \PhiEnt{X} = \E{\Psi(X)} - \Psi(\E{X}) ,
    $$ 
    for $X\geq 0$.
    For any funcion $f$ and distribution $\mu$, define 
    $$
    \PhiEnt[\mu]{f} = \E[\mu]{\Psi(f)} - \Psi(\E[\mu]{f}) ,
    $$ 
\end{definition}
\begin{lemma}[Law of Total Entropy]
    $$
        \PhiEnt{Y} = \E{\PhiEnt {Y\mid X}} + \PhiEnt{\E{Y\mid X}} . 
    $$
\end{lemma}
Suppose we consider the distribution $\pi_k$ where $\supp(\pi_k) \subseteq \binom{[n]}{k}$.
\begin{definition}[Contraction of $\Psi$-entropy]
    We say $\pi_k$ is $C$-contraction of $\Psi$-entropy with respect to $k$ and $r$ if for any $f: \binom{[n]}k\mapsto [0,1]$,
    $$
        \E[S\sim\pi_r]{\PhiEnt[\pi_{k-r}^{S}]{f}} \leq C\cdot \PhiEnt[\mu]{f}
    $$
\end{definition}
\begin{definition}[Approximate Tensorization of $\Psi$-entropy]
    We say $\pi_k$ is $C$-ATE of $\Psi$-entropy if
    $$
        \forall f , \PhiEnt[\mu]{f} \leq C\cdot \E[x\sim \pi_k]{\sum_{i=1}^n \PhiEnt[\pi_1^{x_{-i}}]{f}}
    $$
\end{definition}
\noindent Hope: tensorization can imlpy 
$$
    \PhiEnt{f} \lesssim \E[\pi_k]{|\nabla f|}
$$
where $|\nabla f|$ is the $1$-norm of $f$, i.e., $|\nabla f(x)| = \sum_{i=1}^n |\frac{f(x) - f(x^{\oplus i})}{2}|$.

\begin{definition}[$\Psi$-entropy Independence]
    We say $\pi_k$ is $C$-EI of $\Psi$-entropy if for any $f: \binom{[n]}k\mapsto [0,1]$,
    $$
       k\PhiEnt[\{i\}\sim \pi_1]{\E[T\sim \pi_k]{f(T)\mid i\in T}} \leq C\cdot \PhiEnt[\pi_k]{f} . 
    $$
\end{definition}
Generally, by local-to-global arguement, entropic independence always implies approximate tensorization.
\begin{lemma}\label{lem:EI_to_contraction}
    If $\pi_{k-|S|}^S$ satisfies $\Psi$-entropic independence with $C_{EI}(n-|S|,k-|S|)$ for any subset $S$, then for any $0\leq j\leq k$ and $f: \binom{[n]}k\mapsto [0,1]$, we have that
    $$
        \PhiEnt[T\sim \pi_j]{\E[\pi^T_{k-j}]{f}} \leq (1 - \prod_{i=0}^{j-1}(1-\frac{C_{EI}(n-i,k-i)}{k-i}))\PhiEnt[\pi]{f} 
    $$
\end{lemma}
\begin{proof}
    we prove by induction.
    By law of total entropy, we have that
    \begin{align*}
        \PhiEnt[T\sim \pi_{j+1}]{\E[\pi^T_{k-j}]{f}} &= \PhiEnt[S\sim \pi_{j}]{\E[\pi^S_{k-j}]{f}} + \E[S\sim \pi_j]{\PhiEnt[T\sim \pi^S_{j+1}]{\E[\pi^T_{k-j-1}]{f}}}
        \\&\leq \PhiEnt[S\sim \pi_{j}]{\E[\pi^S_{k-j}]{f}} + \frac{C_{EI}(n-j,k-j)}{k-j}\cdot \E[S\sim \pi_j]{\PhiEnt[T\sim \pi^S_{k-j}]{f}}%+ \E[S\sim \pi_j]{\PhiEnt[T\sim \pi^S_{j+1}]{\E[\pi^T]{f}}}
        \\&= (1 - \frac{C_{EI}(n-j,k-j)}{k-j})\PhiEnt[S\sim \pi_{j}]{\E[\pi^S_{k-j}]{f}} + \frac{C_{EI}(n-j,k-j)}{k-j} \PhiEnt[\pi_k]{f}
        \\&\leq (1 - \frac{C_{EI}(n-j,k-j)}{k-j})(1 - \prod_{i=0}^{j-1}(1 - \frac{C_{EI}(n-i,k-i)}{k-i}))\PhiEnt[\pi_{k}]{f} 
        \\&+ \frac{C_{EI}(n-j,k-j)}{k-j} \PhiEnt[\pi_k]{f}
        \\&= (1 - \prod_{i=0}^{j}(1 - \frac{C_{EI}(n-i,k-i)}{k-i})) \PhiEnt[\pi_k]{f}\qedhere
    \end{align*}
\end{proof}
It is sufficient to build the $\Psi$-entropic independence.

\subsection{Uniform Distribution}
Consdier proving for the uniform distribution on $\pi_k = Uni(\binom{[n]}{k})$.
Take $\Psi = - I = - \varphi(\Phi^{-1})$, where $I$ is the Gaussian isopermetric profile.

We have the following claim.
\begin{conjecture}
    We have that 
    $$
        C_{EI} = \sup_{\substack{f: \PhiEnt[\pi_k]{f} > 0\\f =\Phi(\lambda x)}} \frac{k\PhiEnt[\{i\}\sim \pi_1]{\E[T\sim \pi_k]{f(T)\mid i\in T}}}{\PhiEnt[\pi_k]{f}} .
    $$
\end{conjecture}
The intuition behind the conjecture:
Let $q_i := \binom{n-1}{k-1}^{-1}\sum_{i\in T} f(T)$.
We have that 
$$
    \frac{C_{EI}}k = \sup_{f: \PhiEnt[\pi_k]{f} > 0} \frac{ \frac 1n\sum_{i=1}^n \Psi(q_i)- \Psi(\frac 1n\sum_{i=1}^n q_i)}{ \E{\Psi(f)} - \Psi(\frac 1n\sum_{i=1}^n q_i)}
    \leq \sup_{\*q = (q_1,\ldots,q_n)} \frac{ \frac 1n\sum_{i=1}^n \Psi(q_i)- \Psi(\frac 1n\sum_{i=1}^n q_i)}{(\inf_{f: \*q}\E{\Psi(f)}) - \Psi(\frac 1n\sum_{i=1}^n q_i)}
$$
The Lagrangian function $\+L$ is as follows for $\inf_{f: \*q}\E{\Psi(f)}$: 
$$
    \+L(x,\lambda) = \binom nk^{-1}\sum_{T} \Psi(x_T) - \sum_{i=1}^n\lambda_i(\binom{n-1}{k-1}^{-1}\sum_{i\in T} x_T - q_i) .
$$
Since $\Psi' = - \varphi'(\Phi^{-1})\frac{1}{\varphi(\Phi^{-1})} = \Phi^{-1}\varphi(\Phi^{-1})\frac{1}{\varphi(\Phi^{-1})} = \Phi^{-1}$, we have that 
$$
    \nabla \+L = 0 \implies
    \begin{cases}
        \binom nk^{-1}\Phi^{-1}(x_T) - \binom{n-1}{k-1}^{-1}\sum_{i\in T}\lambda_i = 0 &,\forall T\in\binom{[n]}k
        \\\binom{n-1}{k-1}^{-1}\sum_{i\in T} x_T - q_i = 0 &, \forall i\in[n]
    \end{cases} .
$$ 
Therefore, if $x$ is a local minimizer of $\inf_{f: \*q}\E{\Psi(f)}$ shows that $f(x) = \Phi(\frac nk\inner{x}{\lambda})$.
\begin{observation}
    % By numerical experiments, we observe that for some $b$ satisfying that $|b|$ is small, 
    % $$
    %    \frac{\PhiEnt[\{i\}\sim \pi_1]{\E[T\sim \pi_k]{\Psi(\inner{T}{\lambda} + b)\mid i\in T}}}{\PhiEnt[T\sim\pi_k]{\Psi(\inner{T}{\lambda} + b)}} \xrightarrow{\eps \to 0^+, a.s.} \frac{n-k}{(n-1)k}
    % $$
    % where $\lambda \sim Uni([-\eps,\eps]^n)$.
    By numerical experiments, we observe that
    $$
       \frac{\PhiEnt[\{i\}\sim \pi_1]{\E[T\sim \pi_k]{\Psi(\inner{T}{\lambda})\mid i\in T}}}{\PhiEnt[T\sim\pi_k]{\Psi(\inner{T}{\lambda})}} \xrightarrow{\eps \to 0^+, a.s.} \frac{n-k}{(n-1)k}
    $$
    where $\lambda \sim Uni([-\eps,\eps]^n)$.
\end{observation}
Therefore, we have the following lemma,
\begin{lemma}
    We have that $C_{EI} \geq \frac{n-k}{n-1}$ and it is conjectured that the equality holds.
\end{lemma}
If the equality holds, then by Lemma \ref{lem:EI_to_contraction}, we have taht for any $j\leq k$,
$$
    \PhiEnt[T\sim \pi_j]{\E[\pi^T_{k-j}]{f}} \leq (1 - \prod_{i=0}^{j-1}(1-\frac{n-k}{(k-i)(n-i-1)}))\PhiEnt[\pi]{f} 
$$
The following shall be true,
\begin{align*}
    1 - \prod_{i=0}^{j-1}(1-\frac{n-k}{(k-i)(n-i-1)})
    &\leq 1 - \prod_{i=0}^{j-1}\frac{k-i-1}{k-i} =  \frac jk.
\end{align*}
And it is almost tight if $n\gg k$.
% since $\frac{1}{x} \leq \ln x - \ln(x-1)$ .

Suppose $k \ll n$. In conclusion, the uniform distribution may satisfy $(\frac jk)$-contraction of $(-I)$-entropy with respect to $k$ and $j$ (if the conjecture is right).

Proof idea: fix $v\in \R^n$ and $f_\lambda(x) := \Phi(\lambda\inner{x}{v})$.
Let $A(\lambda) = \frac 1n\sum_{i=1}^n \Psi(q_i)- \Psi(\frac 1n\sum_{i=1}^n q_i), B(\lambda) =  \E{\Psi(f_\lambda)} - \Psi(\frac 1n\sum_{i=1}^n q_i)$ where $q_i  = \binom{n-1}{k-1}^{-1}\sum_{i\in T} f_\lambda(T)$.
Let $g(\lambda) = \frac{A(\lambda)}{B(\lambda)}$.
If we can show that $g$ is increasing with respect to $\lambda > 0$ (which should hold by numerical experiments), then it is sufficient to calculate $\lim_{\lambda\to 0^+} g(\lambda)$, which should converge to $\frac{n-k}{(n-1)k}$.

$\Psi(\frac 12 + \delta) = \Psi(\frac 12) + \frac{\Psi''(\frac 12)}{2}\delta^2 + o(\delta)$.
We have that
$$
    A(\lambda) = \frac{\Psi''(1/2)}2(\frac 1n\sum_{i=1}^n \delta_i^2 - \overline \delta^2) + o(\lambda^2) = \frac{\sqrt{2\pi}}2(\frac 1n\sum_{i=1}^n \delta_i^2 - \overline \delta^2) + o(\lambda^2),
$$
where $\delta_i := q_i - \frac 12 = O(\delta)$, $\overline \delta = \frac 1n\sum_{i=1}^n \delta_i$.
Similarly,
$$
    B(\lambda) = \frac{\lambda^2}{2\sqrt{2\pi}}(\E{\inner{v}{T}^2} - \E{\inner{v}{T}}^2) + o(\lambda^2) = \frac{\lambda^2}{2\sqrt{2\pi}}\Var{\inner{v}{T}}  + o(\lambda^2).
$$
$\delta_i = \E{\frac 12 + \frac{\lambda \inner{v}{T}}{\sqrt{2\pi}} + o(\lambda)\mid i\in T} - \frac 12 = \frac {\lambda}{\sqrt{2\pi}}\E{\inner{v}{T}\mid i\in T} + o(\lambda)$.
Let $\*m_i = \E{\inner{v}{T}\mid i\in T}$ and $\*m = \E{\inner{v}{T}}$, we have that $A(\lambda) = \frac{\lambda^2}{2\sqrt{2\pi}}\Var{\*m_i}$.
Therefore,
\begin{align*}
    \lim_{\lambda \to 0^+} \frac{A(\lambda)}{B(\lambda)} &= \frac{\Var{\*m_i}}{\Var{\inner{v}{T}}} = \frac{\Var[i]{\E{\inner{v}{T}\mid i\in T}}}{\Var{\inner{v}{T}}}
    = \frac{\Var[i]{\frac{n-k}{n-1}v_i + \frac{k-1}{n-1}\sum v_j}}{\Var{\inner{v}{T}}}
    \\&= \frac{(n-k)^2\Var[i]{v_i}}{(n-1)^2\Var{\inner{v}{T}}} = \frac{(n-k)^2(n-1)\Var[i]{v_i}}{(n-1)^2k(n-k)\Var[i]{v_i}} = \frac{n-k}{k(n-1)} ,
\end{align*}
where $\Var{\inner{v}{T}} = \frac{k(n-k)}{n-1}\Var[i]{v_i}$.
% Since $A(0) = B(0) = 0$, we have that
% \begin{align*}
%     A'(\lambda) &= \frac 1n\sum_{i=1}^n \Phi^{-1}(q_i)q_i' - \frac 1n\Phi^{-1}(\frac 1n\sum_{i=1}^n q_i)\sum_{i=1}^n q_i' 
%     \\&= \frac {1}{k\binom nk}\left(\sum_{i=1}^n \sum_{i\in T} \Phi^{-1}(q_i)f_\lambda'(T) - \Phi^{-1}(\frac 1n\sum_{i=1}^n q_i)\sum_{T} f_\lambda'(T)\right) 
%     \\&= \frac {1}{k\binom nk}\sum_{T}f_\lambda'(T)\left(\sum_{i\in T} \Phi^{-1}(q_i) - \Phi^{-1}(\frac 1n\sum_{i=1}^n q_i)\right) ;
% \end{align*}
% \begin{align*}
%     B'(\lambda) &= \frac 1{\binom nk}\sum_T \Phi^{-1}(f_\lambda(T))f_\lambda'(T) - \frac 1n\Phi^{-1}(\frac 1n\sum_{i=1}^n q_i)\sum_{i=1}^n q_i' 
%     \\&= \frac 1{\binom nk}\sum_T \Phi^{-1}(f_\lambda(T))f_\lambda'(T) - \frac 1{k\binom nk}\Phi^{-1}(\frac 1n\sum_{i=1}^n q_i)\sum_{T} f_\lambda'(T)
%     \\&= \frac 1{k\binom nk}\sum_T f_\lambda'(T) \left(k\lambda \inner{T}{v}- \Phi^{-1}(\frac 1n\sum_{i=1}^n q_i)\right) .
% \end{align*}
% Since $\Phi^{-1}(\E{\Phi(\lambda f)}) = \lambda \E{f} + \frac {\lambda^3}6(\E{f}^3 - \E{f^3}) + o(\lambda^4)$,
% $$
%     \sum_{i\in T}\Phi^{-1}(q_i) = \lambda \sum_{i\in T} \E[S\sim \mu(\cdot\mid i\in)]{\inner{S}{v}}  + o(\lambda) = \lambda \E[S\sim \mu]{\inner{S}{v}|S\cap T|} .
% $$
% Want to compute 
% $$
%     \sup_{v} \frac{\sum_T \inner{T}{v}(\E{\inner{S}{v}\cdot|S\cap T|} - \*\mu(\*v))}{\sum_T k\inner{T}{v}^2 - \inner{T}{v}\*\mu(\*v)} ,
% $$
% where $\*\mu(\*v) := \E[S\sim \mu]{\inner{S}{v}}$.


\newpage
\section{Boolean Hypercube Isoperimetric Inequality}
\subsection{Isoperimetric Inequality}
See \cite{durcik2026sharp} for more details.
Let $\Omega = \{0,1\}^n$ be the Boolean hypercube.
Define $h_A(x) := |N(x) \cap (\Omega\setminus A)|$ be the number of edges that go out of $A$ from $x$ for $x\in A$.
\begin{theorem}[Therorem 1 in \cite{durcik2026sharp}]\label{thm:durcik}
    For any $A\subseteq \Omega$ with $|A|\leq \frac 12$,
    we have that 
    $$
        \E[\mu]{\sqrt{h_A}} \geq |A|\sqrt{\log_2(1/|A|)} ,
    $$
    where $\mu$ is the uniform distribution on $\Omega$ and $|A| := |A|/2^n$.
\end{theorem}
\begin{corollary}
    Theorem \ref{thm:durcik} implies that for any boolean function $f: \Omega\mapsto \{0,1\}$, $\norm{\nabla f}_1 \geq \norm{f - \E{f}}_1$, where $\norm{\nabla f}_1 = \E[\mu]{\sqrt{\sum_{i\in[n]}(\frac 12(f(x) - f(x^{\oplus i})))^2}}$.
    And by Cauchy-Schwarz inequality, we have that 
    $$
        \E{h_A} \geq \frac{|A|^2}{|\partial A|}\log_2(1/|A|) .
    $$
\end{corollary}
Proof Idea: find a Bellman function $B: [0,1]\mapsto [0,\infty)$ which satisfies that 
$$
    B(0) = 0, B(1) = 1, max(G_1,G_2)[B](x,y) \geq 0 \text{ for }x,y\in[0,1],
$$
$$
    \text{where }G_1[B](x,y) = \sqrt{(y-x)^2 + B(y)^2} + B(x) - 2B(\frac{x+y}{2}) ,
$$
$$
   G_2[B](x,y) = y-x +(\sqrt 2 - 1)B(y) + B(x) - 2B(\frac{x+y}{2}) .
$$
Then $\E{\sqrt{h_A}} \geq B(|A|)$ for any $A\subset \Omega$.
In \cite{durcik2026sharp}, the authors take
$$
    B(x) = x\sqrt{\log_2(1/x)}\ind[x\leq \frac 14] + \text{some cubic func. for } x\in (\frac 14,\frac 12) +  \frac 12 I(\frac 1{2w})^{-1}I(\frac{1-x}w)\ind[x\geq \frac 12]
$$
to finish the proof.
\subsection{KKL Theorem}
Define $\+I_i(f) := \frac 12\E[\mu]{(f(x) - f(x^{\oplus i}))^2} = \frac 12\Pr{f(x) \neq f(x^{\oplus i})}$ be the influence of the $i$-th coordinate.
Let $\+M(f) := \max_{i\in[n]}\+I_i(f)$ and $\+E(f) := \frac 1n\sum_{i\in[n]}\+I_i(f)$.
\begin{theorem}[KKL]
    For any boolean function $f: \Omega\mapsto \{0,1\}$, we have that
    $$
        \+E(f) \geq\Omega(1)\frac{\log(1/\+M(f))}{n}\Var{f} ,
    $$
    implying $\+M(f) \geq \Omega(\frac {\log n}n)\Var{f}$.
\end{theorem}
\subsection{FKN Theorem}
\begin{theorem}[FKN Theorem]
    Suppose $f: \Omega\mapsto \{0,1\}$ is $\eps$-close to an affine function.
    Then $f$ is $O(\eps)$-close to one of the functions $\{0,1,x_1,1-x_1,\ldots,x_n,1-x_n\}$.
    $f,g$ are $\eps$-close if $\norm{f-g}^2_\mu \leq \eps$, i.e. $\Pr[\mu]{f\neq g}\leq \eps$.

    Specifically, if $\eps = 0$, $f\in \{0,1\}$ or $f\in \{x_i,1-x_i\}$ for some $i\in[n]$.
\end{theorem}
\subsection{Gross-Eldan Inequality}
See \cite{ivanisvili2025eldan} for more details.
\begin{theorem}[Therorem 3 in \cite{ivanisvili2025eldan}]
    There exists a constant $C$ such that for any boolean function $f: \{\pm 1\}^n\mapsto \{\pm 1\}$, we have that
    $$
        \E[\mu]{\norm{\nabla f}_1}\geq C \Var{f} \sqrt{\log\left(1 + \frac{e}{{\color{red} 4}\sum_{i=1}^n \+I_i(f)^2}\right)} .
    $$
\end{theorem}
Proof idea: see \href{https://www.overleaf.com/project/68c323dbe053030b143776e0}{overleaf notes}.
Apply the semigroup $H_t = (e^{-t} I  + (1-e^{-t})\E{})^{\otimes n}$ (i.e., resample w.p. $1 - e^{-t}$ and do nothing w.p. $e^{-t}$).
Prove the hypercontractivity, commutativity ($\E{H_t()} = H_t(\E{})$) and Bobkov inequality on discrete hypercube.
\section{Slices of the Boolean Hypercube}
\subsection{Isoperimetric Inequality} 
What we want to show is that for some constant $C$,
$$
    \E[\mu]{\sqrt {h_S}} \geq C|S|(1-|S|) .
$$
\subsection{KKL Theorem}
See \cite{o2013kkl} for more details.

The definition of influence is as follows: for $1\leq u < v\leq n$,
$$
    \+I_{u,v} = \frac 12\E[\mu]{(f(x) - f(x^{u,v}))^2} = \frac 12\Pr{f(x) \neq f(x^{u,v})} ,
$$
where $x^{u,v}$ denotes the vector obtained by swapping the $u$-th and $v$-th coordinates of $x$.
We have the following similar KKL theorem for the slice of the Boolean hypercube.
\begin{theorem}[Theorem 1.1 in \cite{o2013kkl}]\label{thm:kkl_slice}
    Suppose the uniform distribution over $\binom{[n]}k$ has log-Sobolev constant $\rho$, we have that for any boolean function $f: \binom{[n]}{k}\mapsto \{0,1\}$, 
    $$
        \+E(f) \geq \Omega(1)\rho\log(1/\+M(f))\Var{f} ,
    $$
    implying $M(f)\geq \Omega(1)\rho \log(1/\rho)\Var{f}$.
    Specifically, if $\Omega(1) \leq k/n\leq 1-\Omega(1)$, we can replace $\rho$ with $\frac 1n$, since $\rho = \Theta\left(\displaystyle\frac 1{n\log \frac{n(n-1)}{2k(n-k)}}\right) = \Theta\left(\frac 1n\right)$.
\end{theorem}
Proof idea: define generator $L = \binom{n}{2}^{-1}\sum_{a<b} L_{a,b}$ such that $L_{a,b}f(x) = f(x) - f(x^{a,b})$ (i.e., $L = I - K$ where $K$ is down-up walk on the slice with some laziness).
Then $H_t = e^{tL}$ satisfies that $H_t f(x) = \E[m\sim \-{Poisson}(t), y\sim K^{m}(x,\cdot)]{f(y)}$.
For uniform distribution, we have bounds for spectral gap and log-Sobolev constant:
$$
    \Theta\left(\displaystyle\frac 1{n\log \frac{n(n-1)}{2k(n-k)}}\right) = \rho \leq \lambda_1 = \frac 2{n-1} 
$$
hypercontractivity:
$$
    \forall 1\leq p\leq q\text{ with } \frac{q-1}{p-1} \leq \exp(2\rho t) \text{ and } f\in L^2(\mu): \norm{H_t f}_q \leq \norm{f}_p .
$$
and commutativity:
$$
    L L_{a,b} = L_{a,b} L \text{ and }\forall t\geq 0, H_t L_{a,b} = L_{a,b} H_t .
$$
Calculation shows Theorem \ref{thm:kkl_slice}.
\subsection{FKN Theorem}
See \cite{filmus2014friedgut} for more details.
\begin{theorem}[Theorem 3.1 in \cite{filmus2014friedgut}]
    Suppose a boolean function $f: \binom{[n]}{k}\mapsto \{0,1\}$ is $\eps$-close to an affine function, where $2\leq k\leq n-2$.
    Let $p:=\min\{k/n,1-k/n\}$.
    Then either $f$ or $1 - f$ is $O(\eps)$-close to $\max_{i\in S} x_i$ or to $\min_{i\in S} x_i$ for some $S\subseteq [n]$ of size at most $\max\{1,O(\sqrt{\eps})/p\}$.

    Specifically, if $\eps = 0$, $f\in \{0,1\}$ or $f\in \{x_i,1-x_i\}$ for some $i\in[n]$.
\end{theorem}


% \subsection{Notation Translation}
% For $f\in C_k = (X_k,\R)$, the influence (only one layer of down-up walk)
% \subsection{Verification}
% The KKL theorem is that
% $$
%     \+E(f) \geq \Omega(1)\frac{\log(1/\+M[f])}{n}\Var{f} .
% $$
$(\eps,i)$-global if 
$$
    \forall \tau \in X(i): \E[X_\tau]{f_\tau} \leq \E{f} + \eps.
$$
% If $\eps$ is large, it means that the density must concentrate on some subset of the simplicial complex, meaning that some influence must be large.
% Use Garland’s Lemma to do induction.
\subsection{Main Results}
$\gamma_i = \max_{\tau\in X(i)}\{\lambda_2(A_\tau)\}, \gamma_i^{(-)} = \min_{\tau\in X(i)}\{\lambda_{\min}(A_\tau)\}$.
\begin{theorem}[Theorem 1.1 and Theorem 1.2 in \cite{hopkins2025local}]
If $f$ is $(\eps,i)$-global boolean function, Then
$$
    \Phi(f) \geq \frac{1-\E{f}}{k-i}\prod_{j=i}^{k-2}(1-\gamma_j) - c_{k,i,\gamma} \eps ,
$$
where $c_{k,i,\gamma} \leq (k-i+1)\left(\frac 1{k-i}\prod_{j=i}^{k-2}(1-\gamma_j)-\frac 1{k-i+1}\prod_{j=i-1}^{k-2}(1-\gamma_j)\right)\left(1+(k-i)\gamma_{i-1}^{(-)}\right)^{-1}$ and $\Phi$ is the conductance of the down-up walk on $C_k$.

And for $0<i<k$, there exists $\{X_k\}$ and $(0,i)$-global function $\{f\}$ such that 
$$
    \Phi(f)\leq (1 - \E{f})\frac {i+1}k\prod_{j=i}^{k-2}(1-\gamma_j),
$$
where, $\Phi(f) = 1-\frac{\inner{N_k^1f}{f}}{\inner{f}{f}}$.
\end{theorem}
\begin{example}
    Let $X(k) = \binom{[n]}k$ and $\pi_k$ be the uniform distribution on $X(k)$.
    We have that
    $$
        \tau\in X(i),A_\tau = \frac 1{n-|\tau|-1}(J-I)\implies \gamma_i = \gamma_i^{(-)} = -\frac 1{n-i-1} .
    $$
    $$
        c_{k,i,\gamma} \leq \frac{n-i}{(n-k+1)(k-i)}.
    $$
    Note that under worst case,
    $$
        \eps = \max_{\tau}\{\E[X_\tau]{f_\tau} - \E{f}\} = \max_\tau \frac{|S\cap X_\tau|}{\binom {n-i}{k-i}} - \frac{|S|}{\binom {n}{k}} = 1 - \frac{\binom {n-i}{k-i}}{\binom {n}{k}} = 1 - \frac{k^{\underline i}}{n^{\underline i}}.
    $$
    Therefore,
    \begin{align*}
        \Phi(f) &\geq \frac{1-\E{f}}{k-i}\prod_{j=i}^{k-2}(1-\gamma_j) - c_{k,i,\gamma} \eps
        \\&= \frac{1-\E{f}}{k-i}\frac{n-i}{n-k+1} - \frac{n-i}{(n-k+1)(k-i)}\eps
        \\&= \frac{n-i}{(n-k+1)(k-i)}(1 - \E{f} - \eps)
        \\&= \frac{n-i}{(n-k+1)(k-i)}(1 - \max_\tau \E[X_\tau]{f_\tau})
        \\&= \frac{n-i}{(n-k+1)(k-i)}(1 - \max_{\tau:|\tau| = i} \frac{|S\cap X_\tau|}{\binom {n-i}{k-i}}) .
    \end{align*}
    where $S = \supp(f)$.
\end{example}

\bibliographystyle{alpha}
\bibliography{ref.bib}
\end{document}